\begin{figure}[!ht]
    \centering
    \resizebox{\textwidth}{!}{%
    \begin{tikzpicture}[
        font=\scriptsize,
        node distance = 9mm,
        stage/.style   = {rectangle, draw, rounded corners, fill=blue!8,
                          text width=2.7cm, align=center, minimum height=1.5cm},
        note/.style    = {text width=2.7cm, align=center, font=\tiny, text=black!65},
        arr/.style     = {-{Stealth[length=2mm]}, thick},
    ]

    \node[stage] (capture)  {\textbf{1. Capture}\\Synchronized stereo pair};
    \node[stage, right=of capture]  (seg)      {\textbf{2. Segmentation}\\Image $\rightarrow$ binary cable mask};
    \node[stage, right=of seg]      (endpoints){\textbf{3. Endpoints}\\Mask $\rightarrow$ the two loose ends};
    \node[stage, right=of endpoints](track)    {\textbf{4. Path tracing}\\Ends $\rightarrow$ ordered points along the thread};
    \node[stage, right=of track]    (match)    {\textbf{5. Stereo matching}\\Left/right path correspondence};
    \node[stage, right=of match]    (metrics)  {\textbf{6. Triangulation}\\3D coordinates \& metrics};

    \draw[arr] (capture) -- (seg);
    \draw[arr] (seg) -- (endpoints);
    \draw[arr] (endpoints) -- (track);
    \draw[arr] (track) -- (match);
    \draw[arr] (match) -- (metrics);

    \node[note, below=4mm of seg]   {HSV threshold (classical) \emph{or} Random Forest classifier (AI, implemented)};
    \node[note, below=4mm of track] {SmartWireTracker search (classical, implemented) \emph{or} learned scorer (AI, proposed)};

    \end{tikzpicture}}
    \caption[Implementation-level view of the Phase II reconstruction pipeline]{Implementation-level view of the Phase II (geometry-driven) reconstruction pipeline. The six boxes are the processing stages, each labeled with the transformation it performs. Where a stage currently offers more than one implementation -- segmentation and path tracing -- the classical method already implemented and the AI-assisted alternative are noted underneath; all other stages have a single, classical implementation today. Algorithm-level detail for each stage is given in the surrounding text rather than in the diagram.}
    \label{fig:pipeline-implementation}
\end{figure}
